\documentclass[11pt]{article}
\usepackage[margin=1in]{geometry}
\usepackage{amsmath,amssymb,amsthm}
\usepackage{booktabs}
\usepackage{array}
\usepackage{xcolor}
\usepackage{hyperref}
\hypersetup{colorlinks=true,linkcolor=blue!50!black,urlcolor=blue!50!black,citecolor=blue!50!black}
\usepackage{fancyhdr}
\usepackage{longtable}
\usepackage[protrusion=true,expansion=false]{microtype}

\pagestyle{fancy}
\fancyhf{}
\fancyhead[L]{\small The Geometry Engine}
\fancyhead[R]{\small The Geometry Engine}
\fancyfoot[C]{\thepage}

\newtheorem{theorem}{Theorem}
\newtheorem{proposition}{Proposition}
\newtheorem{observation}{Observation}
\newtheorem{definition}{Definition}
\newtheorem{principle}{Principle}

\title{\textbf{The Geometry Engine:\\
Multiplicative Structure as Shape, the Additive Seam,\\
and the Two Faces of the Riemann Zeta Function}}
\author{B.\ Wallace \quad (TensorRent)\\
\small with computational verification by Claude (Anthropic)}
\date{\small Companion to Papers~a--f \ ---\ SIP License v1.1 \ ---\ Seed \texttt{20260423}}

\begin{document}
\maketitle

\begin{abstract}
We develop a representation in which a number's \emph{multiplicative} structure is a geometric object --- a
box whose sides are primes and whose dimensions are exponents --- and operations that conserve volume are
answered directly from the shape without computing the value. A 91-line reference engine implements this;
all twelve of its operations are verified exact against integer ground truth. Benchmarked head-to-head
against standard methods across ten ontologies of applied mathematics and computer science, the engine wins
decisively where the standard method is forced to flatten rich multiplicative structure --- multiplicative
functions ($81\times$), the prime structure of binomial coefficients without expansion ($13\times$),
exact rationals, canonical fingerprinting --- and loses, flatly reported, where the query needs only partial
structure (smoothness, $106\times$ slower), where the standard operation is already trivial (\textsc{gcd}),
or where the task is set-generation (sieving, $73\times$ slower); it has no advantage at the factoring entry
wall, by design. The representation is bounded by one seam: \emph{addition has no shape}. We confirm, in three
independent settings (the Riemann zeros, Goldbach decompositions, and prime gaps), a single law --- multiplication
governs the \emph{density} of the additive world while individual additive structure stays generic, the bridge
being statistical and never an address --- and we report honestly that the coupling is strong for first-order
counts ($\sim\!0.99$) and only moderate for second-order differences ($\sim\!0.6$), including a sub-claim that
failed its pre-registered test. The engine generalizes to any unique-factorization domain, and to analytic
functions with poles as their ``primes.'' Finally, on the Riemann zeta function the engine's two faces --- prime
boxes (the Euler product) and zero/pole geometry (the Hadamard product) --- are the same object joined by the
explicit formula, which we verify in both directions on the real first $10^5$ zeros. Throughout, classical
results are labeled as verification rather than novelty, negative results are retained as first-class outputs,
and the applied target is overhead reduction, for which a battle-tested known law is an asset.
\end{abstract}

\section{Orientation and honesty protocol}

This paper records one working session conducted under a fixed discipline: every claim passes a
pre-registered computational test before it is stated; results are labeled by \emph{how} they were verified;
negative results are kept as first-class outputs with the killing computation; and scope boundaries are
stated explicitly. The applied goal is reduction of computational overhead, not mathematical novelty ---
consequently a result already established in the literature is an asset (battle-tested, buildable), not a
demerit. We use a four-tier verification ledger: \textbf{(T1)} machine-verified against ground truth;
\textbf{(T2)} proven or forced; \textbf{(T3)} numerically verified, not theorem-level; \textbf{(T4)}
explicitly falsified.

\section{The representation: numbers as shapes}

\begin{definition}[Shape]
A \emph{shape} is a box, encoded as a finite map $\{\,p_i : e_i\,\}$ from prime sides $p_i$ to dimension
exponents $e_i \ge 1$. Its \emph{volume} is the integer $\prod_i p_i^{e_i}$. The empty box is the unit $1$.
Zero and negatives are not boxes.
\end{definition}

\begin{principle}[Conservation]
Every operation on shapes is admissible if and only if it conserves volume. This single invariant is the
correctness criterion: a reshaping is exact when the volume it produces equals the value the corresponding
arithmetic operation would produce.
\end{principle}

The content of the representation is that the \emph{multiplicative} questions about a number are
\emph{geometric} questions about its box, answerable without ever computing the volume.

\begin{theorem}[Shape-native operations, T1]
The following operations are computable from the shape alone and were verified exact against integer ground
truth at the indicated scale (seed \texttt{20260423}):
\begin{center}
\begin{tabular}{lll}
\toprule
operation & geometry & verification \\
\midrule
multiply & merge boxes (exponents add) & $2000$ pairs \\
power & scale all dimensions by $k$ & samples $\times\, k\in\{0,1,2,3\}$ \\
divisibility & box containment & $2000$ pairs \\
exact division & remove the nested box & $2000$ pairs \\
\textsc{gcd} / \textsc{lcm} & common / enclosing box (min / max dims) & $2000$ pairs \\
primality & atomic box: one side, unit dimension & all $[2,2\!\times\!10^4)$ \\
perfect power & dimensional symmetry: $\gcd$ of dims $\ge 2$ & all $[2,5000)$ \\
factorization & the box \emph{is} the decomposition & $500$ samples \\
\bottomrule
\end{tabular}
\end{center}
\end{theorem}

The reference implementation is $91$ lines; the harness reports $12/12$ exact, including the round-trip
$\mathrm{volume}(\mathrm{shape}(n))=n$ and the glyph rendering (sides named by prime index, the
``built-once, shared'' dictionary).

\subsection{The boundary: addition has no shape}

\begin{proposition}[The additive seam, T1]
Addition admits no box-merge: the sum of two boxes is in general not a box, and its shape is unrelated to
the shapes of the summands. Concretely the $360$-box and the $84$-box sum to $444=2^2\cdot 3\cdot 37$, a box
sharing no structure with either input.
\end{proposition}

In the engine this is handled by an explicit, named \texttt{add\_exit} that computes the value, refactors,
and re-enters geometry, emitting a \texttt{GEOMETRY\_EXIT} flag so that every additive step is visible as a
departure from the geometric world. Addition is never hidden inside a geometric costume.

\section{The advantage map (measured)}

We raced the engine against the standard method for each task, reporting wins and losses flatly.

\begin{center}
\begin{tabular}{llrl}
\toprule
ontology & task & factor & outcome \\
\midrule
multiplicative functions & $\varphi,\sigma,\mu,d(n)$ on held shapes & $81\times$ & \textbf{win} \\
combinatorics & prime structure of $\binom{n}{k}$, no expansion & $13\times$ & \textbf{win} \\
combinatorics & $\binom{5000}{2500}$ value via Legendre & $2.2\times$ & win \\
exact rationals & product of fractions, auto-reduced & exact vs.\ lossy & \textbf{win (correctness)} \\
content addressing & canonical path-independent signature & --- & \textbf{win (determinism)} \\
generality & extends to any UFD (polynomials, $\mathbb{Z}[i]$) & --- & \textbf{win (reach)} \\
computer algebra & \textsc{gcd} on pre-factored inputs & $2.9\times$ & loss (Euclid) \\
number theory & coprimality at scale & $2.2\times$ & loss (\textsc{gcd}) \\
number theory & $B$-smoothness detection & $106\times$ & loss (trial division) \\
cryptography & factor an arbitrary semiprime & --- & no advantage (by design) \\
prime generation & generate primes below $N$ & $73\times$ & loss (sieve) \\
\bottomrule
\end{tabular}
\end{center}

\begin{observation}[The predictive law]
The advantage is proportional to the multiplicative structure the standard method is forced to flatten,
minus the cost of building the shape. The engine wins when the shape is already held (or is cheap because
the number is structured) and the query is multiplicative; it loses when the query needs only partial
structure (smoothness needs only small primes), when the standard operation is already trivial
(\textsc{gcd} on machine integers), or when the task is set-generation (sieving shares work across a range
that a per-number tool cannot). It has no advantage where building the shape is itself the hard problem
(factoring) --- necessarily so, since otherwise public-key cryptography would fail.
\end{observation}

The combinatorial win is the sharpest: the question ``what power of $7$ divides $\binom{5000}{2500}$'' is
answered in $0.34$\,ms from the shape (via Legendre's exponent formula) versus $4.4$\,ms to build the
$\sim\!16000$-digit integer and factor it --- a $13\times$ gain that never forms the giant number, the direct
payoff of holding an object in its native dimension rather than its flattened projection.

\section{The multiplicative--additive seam}

\begin{principle}[Seam law, T3 across three instances]
Multiplication governs the \emph{density} of the additive world; addition fills it generically; the bridge
between the two is statistical, never an address.
\end{principle}

\begin{center}
\begin{tabular}{llll}
\toprule
instance & individual structure & density governed by & coupling \\
\midrule
Riemann zeros & spacings generic (GUE) & positions $\leftarrow$ primes & $+122\sigma$ (strong) \\
Goldbach & decomposition a cloud & count $\leftarrow$ Hardy--Littlewood & $0.9952$ (strong) \\
prime gaps & gaps generic (Poisson) & density $\leftarrow$ small-prime richness & $0.59$ (moderate) \\
\bottomrule
\end{tabular}
\end{center}

\begin{observation}[Honest variation, including a falsified sub-claim, T4]
The coupling strength varies with the order of the additive object: strong for first-order counts
($\sim\!0.99$), moderate for second-order differences ($\sim\!0.6$). The ``jumping champions are primorials''
sub-claim \emph{failed} its pre-registered test at $N=3\times10^6$ (the empirical champions are $6,12,2,4$;
the primorial ladder dominates only near $\sim\!10^{35}$) and is recorded as falsified at this scale.
\end{observation}

\begin{observation}[Applied corollary, mixed result]
Strong coupling converts to a real computational skip: replacing Goldbach enumeration by the
Hardy--Littlewood density formula measured $223\times$ faster (the formula's shape tracks the truth at
$0.9952$; obtaining single-digit relative error requires properly normalizing the correction constant,
which the in-session demonstration did not). Moderate coupling does \emph{not} automatically yield a usable
predictor: a proposed gap-pruning proxy tested at correlation $-0.029$ and was discarded. High analysis
correlation does not guarantee a cheap exploitable predictor; the specific decision-time proxy must itself
be tested.
\end{observation}

\section{The two-structure map of number}

Every number carries two independent decompositions: its \emph{multiplicative} structure (the unique
prime box, by the Fundamental Theorem of Arithmetic --- representable, a coordinate) and its
\emph{additive} structure (a sum of primes --- whose parts are non-unique, a cloud, not representable, but
whose minimal \emph{depth} is a unique invariant). The depth is $1$, $2$, or $3$: every prime has depth $1$;
every even $\ge 4$ has depth $2$ (Goldbach, verified empirically); every remaining odd has depth $3$ (weak
Goldbach, proven by Helfgott, 2013). Verified over $[2,10^5)$:

\begin{theorem}[The diagonal, T1]
Every prime has additive depth $1$, and only primes have depth $1$. The multiplicative atoms are exactly
the additive atoms: primes are the unique volumes rigid under both reshapings, the single class where
``Form'' (additive) and ``Function'' (multiplicative) coincide. The additively-reachable primes
(those with $p-2$ prime) are precisely the upper members of twin-prime pairs.
\end{theorem}

\section{The analysis engine: poles as the primes of a function}

A Taylor \emph{series value} is an additive object and has no box. But the \emph{function} carries genuine
geometry in its singularities, and the engine transfers wholesale with poles in place of primes.

\begin{center}
\begin{tabular}{lll}
\toprule
 & integer engine & analysis engine \\
\midrule
irreducible units & primes & poles \\
shape & $\{p:e\}$ & $\{\text{pole}:\text{order}\}$ \\
free structural questions & divisibility, \textsc{gcd}, $\varphi$ & radius of convergence, best center \\
value (still work) & $\mathrm{volume}()$ & summing the series \\
wall (no finite shape) & generic large numbers & entire / lacunary functions \\
\bottomrule
\end{tabular}
\end{center}

Verified: the radius of convergence equals the distance from the expansion center to the nearest pole
($1/(1-x)\Rightarrow R=1$, matching the coefficient ratio test). For the real function $1/(1+x^2)$, whose
real Taylor series converges only for $|x|<1$ with no reason visible on the real line, the geometry supplies
the reason --- poles at $\pm i$, distance $1$ --- which summation can only measure. The optimal expansion
center is chosen by pure geometry (farthest from all poles), summing no series at any candidate.

\section{Loop closure: the two faces are one engine}

\begin{theorem}[Two faces of zeta, T1 both directions]
The integer engine (prime boxes) and the analysis engine (pole/zero geometry) are one engine seen through
the two products of the Riemann zeta function: the Euler product $\prod_p (1-p^{-s})^{-1}$ (primes as
multiplicative boxes) and the Hadamard product (zeros as analytic pole geometry). The explicit formula is
the dictionary between them, verified in both directions on the real first $10^5$ zeros:
\begin{itemize}\setlength\itemsep{1pt}
\item zeros' positions reconstructed from primes: $+122\sigma$ over an amplitude-matched random control;
\item the prime staircase $\psi(x)$ reconstructed from zeros: agreement at every tested $x$
($7.83/7.82$, $49.5/49.5$, $94.0/94.2$), tightening as more zeros are included.
\end{itemize}
\end{theorem}

Thus the two threads --- geometric number representation and the Riemann zeros --- are the same structure.
At the top level: primes are the multiplicative Form, zeros are the analytic Form, and zeta is the Function
carrying both.

\section{Scope and what is owned}

The mathematics underlying \S6--\S7 is classical: the Fundamental Theorem of Arithmetic, Legendre's
exponent formula, the Hardy--Littlewood conjecture, the Hadamard product, and the Riemann--von Mangoldt
explicit formula. The verifications confirm that the engine's logic \emph{extends} to those domains; they
are not new theorems, and for the applied goal of overhead reduction their established status is an asset.
What is genuinely contributed here is structural: a single shape-engine architecture --- irreducibles yield
a representable shape, structural questions are free, the value remains additive work, and infinite structure
breaks the finite box --- is shown to be the \emph{same} architecture on the prime side and the zero side,
and the advantage map and seam law are measured rather than asserted. Reconstructions are finite-$N$ and
approximate (Gibbs ringing, finite zero count, a $\sim\!5\%$ staircase band). A Rust/OmniForge transcription
is pending: the verification environment lacked a Rust toolchain, so all artifacts are a Python reference
implementation plus a locked specification, and no unverified Rust was produced. Falsified claims are
retained: Mersenne-as-prime, single-deletion-as-primality, geometry-beats-division, the gap-pruning
heuristic ($-0.029$), and the primorial-champions sub-claim.

\vspace{1em}
\noindent\rule{\linewidth}{0.4pt}\\
\small Companion to \textit{Colour from Gravity} through \textit{Three Layer Decomposition} (Papers~a--f). Reference engine, harness, and specification accompany this paper.
SIP License v1.1. Seed \texttt{20260423}.

\end{document}
